\subsection{How fast does the relaxation approach the ground?}
\label{subsec:relaxation-rate}
The proof of Proposition~\ref{prop:two-moment-deconditioning} used that every state relaxes to $P_y$. Here we give the rate and the generator. Write $\|X\|_1=\Tr|X|$ for the trace norm, and for an appearance state $\sigma$ let
\begin{equation}
 \delta_y(\sigma)=\Tr(Q_y\sigma)=1-\langle y,\sigma y\rangle
 \label{eq:relaxation-deviation}
\end{equation}
be its weight outside the ground. The complementary weight $\langle y,\sigma y\rangle$ is the probability that $\sigma$ passes the ground test $P_y$.

\begin{proposition}[Distance from a pure state]
\label{prop:relaxation-pure-distance}
For a density operator $\sigma$ and a unit vector $y$,
\begin{equation}
 2\,\delta_y(\sigma)\;\le\;\|\sigma-P_y\|_1\;\le\;2\sqrt{\delta_y(\sigma)} .
 \label{eq:relaxation-pure-distance}
\end{equation}
The upper bound is an equality when $\sigma$ is pure, and the lower bound when $\sigma$ commutes with $P_y$. In particular $\delta_y(\sigma)=0$ exactly when $\sigma=P_y$.
\end{proposition}
\begin{proof}
The operator $Q_y-P_y$ has norm one, so $\|\sigma-P_y\|_1\ge\Tr\bigl((Q_y-P_y)(\sigma-P_y)\bigr)=2\delta_y(\sigma)$. For a unit vector $\phi$, the operator $P_\phi-P_y$ is self-adjoint with trace zero and rank at most two, and $\Tr(P_\phi-P_y)^2=2-2|\langle y,\phi\rangle|^2$. Its nonzero eigenvalues are therefore $\pm\sqrt{1-|\langle y,\phi\rangle|^2}$, and $\|P_\phi-P_y\|_1=2\sqrt{\delta_y(P_\phi)}$. Now write $\sigma=\sum_k p_kP_{\phi_k}$ with orthonormal eigenvectors $\phi_k$ and $\sum_kp_k=1$. The triangle inequality and the concavity of the square root give
\[
 \|\sigma-P_y\|_1\le\sum_kp_k\,2\sqrt{\delta_y(P_{\phi_k})}
 \le2\sqrt{\textstyle\sum_kp_k\,\delta_y(P_{\phi_k})}=2\sqrt{\delta_y(\sigma)} .
\]
If $\sigma$ commutes with $P_y$, then $\sigma-P_y=-\delta_y(\sigma)P_y+Q_y\sigma Q_y$, and the two terms have orthogonal supports and trace norm $\delta_y(\sigma)$ each.
\end{proof}
The upper bound is the inequality of Fuchs and van de Graaf for a pure target \cite{FuchsVandeGraaf1999}. For a density operator, passing the ground test with probability one, having no weight outside the ground and being the state $P_y$ are therefore the same condition.

\begin{proposition}[The relaxation as a semigroup]
\label{prop:relaxation-semigroup}
Let $\gamma>0$ and $s=e^{-\gamma t/2}$. Extend Eq.~\eqref{eq:moment-deconditioning} to every trace-class operator $X$ by $\mathcal D_t^y(X)=L_yXL_y+(1-s^2)\Tr(Q_yX)P_y$.
\begin{enumerate}
\item[(a)] $\mathcal D_{t+r}^y=\mathcal D_t^y\mathcal D_r^y$ for $t,r\ge0$, and $\|\mathcal D_t^y(X)-X\|_1\le4(1-e^{-\gamma t})\|X\|_1$.
\item[(b)] $\mathcal D_t^y=\exp(t\,\mathcal G_y)$, where
\begin{equation}
 \mathcal G_y(X)=\gamma\Bigl(\Tr(Q_yX)\,P_y-\tfrac12\bigl(Q_yX+XQ_y\bigr)\Bigr),
 \qquad \|\mathcal G_y(X)\|_1\le2\gamma\|X\|_1 .
 \label{eq:relaxation-generator}
\end{equation}
\item[(c)] For $t>0$ the trace-class operators fixed by $\mathcal D_t^y$ are the multiples of $P_y$. Thus $P_y$ is fixed, and it is the only fixed state.
\item[(d)] For every state $\sigma$,
\begin{equation}
 \delta_y\bigl(\mathcal D_t^y(\sigma)\bigr)=e^{-\gamma t}\,\delta_y(\sigma).
 \label{eq:relaxation-decay}
\end{equation}
\item[(e)] For every state $\sigma$,
\begin{equation}
 2e^{-\gamma t}\,\delta_y(\sigma)\;\le\;\bigl\|\mathcal D_t^y(\sigma)-P_y\bigr\|_1\;\le\;2e^{-\gamma t/2}\sqrt{\delta_y(\sigma)} .
 \label{eq:relaxation-distance}
\end{equation}
\end{enumerate}
\end{proposition}
\begin{proof}
Write $P=P_y$ and $Q=Q_y$. In the four blocks of $X$ relative to $P$ and $Q$ the map reads
\begin{equation}
 \mathcal D_t^y(X)=PXP+s\,(PXQ+QXP)+s^2\,QXQ+(1-s^2)\Tr(QX)\,P .
 \label{eq:relaxation-blocks}
\end{equation}
A second step with factor $s'$ multiplies the off-diagonal blocks by $ss'$ and the block $QXQ$ by $s^2s'^2$. The coefficient of $\Tr(QX)P$ becomes $(1-s^2)+s^2(1-s'^2)=1-s^2s'^2$. This is the composition law. Subtracting $X$ from Eq.~\eqref{eq:relaxation-blocks} leaves $(s-1)(PXQ+QXP)+(s^2-1)QXQ+(1-s^2)\Tr(QX)P$. Since $\|PXQ+QXP\|_1\le2\|X\|_1$, $\|QXQ\|_1\le\|X\|_1$, $|\Tr(QX)|\le\|X\|_1$ and $1-s\le1-s^2$, its trace norm is at most $4(1-s^2)\|X\|_1$. This proves (a).

The coefficients $s$ and $s^2$ have derivatives $-\gamma/2$ and $-\gamma$ at $t=0$. The derivative of Eq.~\eqref{eq:relaxation-blocks} there is $-\tfrac{\gamma}{2}(PXQ+QXP)-\gamma\,QXQ+\gamma\Tr(QX)P$, which is $\mathcal G_y(X)$ because $QX+XQ=PXQ+QXP+2QXQ$. The three terms of $\mathcal G_y(X)$ in Eq.~\eqref{eq:relaxation-generator} have trace norms at most $\gamma\|X\|_1$, $\tfrac{\gamma}{2}\|X\|_1$ and $\tfrac{\gamma}{2}\|X\|_1$. The difference quotients of $\mathcal D_t^y$ at $t=0$ converge in the operator norm on the trace class, because only the scalar coefficients of the fixed block maps depend on $t$. With the composition law this gives $\tfrac{d}{dt}\mathcal D_t^y=\mathcal G_y\mathcal D_t^y$ and hence (b).

If $\mathcal D_t^y(X)=X$ with $0<s<1$, comparing blocks in Eq.~\eqref{eq:relaxation-blocks} gives $PXQ=QXP=0$ and $QXQ=0$, so $X=PXP$ is a multiple of $P$. Conversely $\mathcal D_t^y(P)=P$. This is (c). Taking the trace of Eq.~\eqref{eq:relaxation-blocks} against $Q$ gives (d). Applying Proposition~\ref{prop:relaxation-pure-distance} to the state $\mathcal D_t^y(\sigma)$ and inserting (d) gives (e).
\end{proof}

Part (a) says more than continuity of $t\mapsto\mathcal D_t^y(X)$ for each $X$: the maps converge to the identity uniformly on bounded sets of the trace class, which is why the generator is bounded. In an orthonormal basis $(e_j)$ of $y^\perp$ it has the standard form of a generator of a semigroup of channels \cite{GKS1976,Lindblad1976}, with no Hamiltonian term,
\begin{equation}
 \mathcal G_y(X)=\sum_j\Bigl(J_jXJ_j^*-\tfrac12\bigl(J_j^*J_jX+XJ_j^*J_j\bigr)\Bigr),
 \qquad J_j=\sqrt\gamma\,\ket{y}\!\bra{e_j}.
 \label{eq:relaxation-jumps}
\end{equation}
Indeed $\sum_jJ_jXJ_j^*=\gamma\Tr(Q_yX)P_y$ and $\sum_jJ_j^*J_j=\gamma Q_y$. Each $J_j$ moves the amplitude along $e_j$ to the ground.

By Eq.~\eqref{eq:relaxation-decay} the weight outside the ground obeys $\tfrac{d}{dt}\delta_y=-\gamma\,\delta_y$. It decreases strictly while it is positive, and it is positive at every finite time if it was positive at the start. The relaxation therefore reaches $P_y$ at a finite time only from $P_y$. Reaching it in one step takes a different channel, such as the reset $\sigma\mapsto P_y$. This is the value $s=0$ of Eq.~\eqref{eq:moment-deconditioning}, which no finite time attains.

The lower bound in Eq.~\eqref{eq:relaxation-distance} is attained: if $\sigma$ commutes with $P_y$, so does $\mathcal D_t^y(\sigma)$, and the distance is $2e^{-\gamma t}\delta_y(\sigma)$. The upper bound is an equality for pure states at $t=0$. At a positive time it is strict unless $\sigma=P_y$. To see this, write $\mathcal D_t^y(\sigma)=c\,\tau+(1-c)P_y$, where $c=1-(1-s^2)\delta_y(\sigma)$ and $\tau=L_y\sigma L_y/c$ is a state with $\delta_y(\tau)=s^2\delta_y(\sigma)/c$. Then $\mathcal D_t^y(\sigma)-P_y=c\,(\tau-P_y)$, and Proposition~\ref{prop:relaxation-pure-distance} applied to $\tau$ gives
\begin{equation}
 \bigl\|\mathcal D_t^y(\sigma)-P_y\bigr\|_1
 \le2s\sqrt{\delta_y(\sigma)}\,\sqrt{1-(1-s^2)\,\delta_y(\sigma)},
 \label{eq:relaxation-sharp-distance}
\end{equation}
with equality when $\sigma$ is pure, because $\tau$ is then pure. For $t>0$ and $\delta_y(\sigma)>0$ the last factor is less than one. It tends to one as $\delta_y(\sigma)\to0$, so at any fixed time the distance for pure inputs close to the ground comes arbitrarily close, in ratio, to the upper bound of Eq.~\eqref{eq:relaxation-distance}. The constant $2$ is therefore optimal, although that bound is not attained at a positive time. The rate $\gamma/2$ is optimal as well, and a single input shows it. For a fixed pure input with $0<\delta_y(\sigma)<1$, the equality in Eq.~\eqref{eq:relaxation-sharp-distance} gives $e^{\gamma t/2}\,\|\mathcal D_t^y(\sigma)-P_y\|_1\to2\sqrt{\delta_y(\sigma)\bigl(1-\delta_y(\sigma)\bigr)}>0$ as $t\to\infty$. The weight outside the ground decays at rate $\gamma$. The coherence between the ground and its complement decays at rate $\gamma/2$, and a state with such coherence approaches the ground at this slower rate.

\subsection{Fixing the ground and attracting to it}
\label{subsec:relaxation-unitary}
Two properties of the relaxation can be asked of any family of channels $\mathcal T_t$ on the appearance space. The family \emph{fixes the ground} if $\mathcal T_t(P_y)=P_y$ for every $t$. It is \emph{ground-attractive} on a set $\mathsf B$ of states if it fixes the ground and $\delta_y(\mathcal T_t(\sigma))\to0$ as $t\to\infty$ for every $\sigma$ in $\mathsf B$. By Eq.~\eqref{eq:relaxation-pure-distance} the second condition is the same as $\mathcal T_t(\sigma)\to P_y$ in trace norm. The set $\mathsf B$ is the basin, and it has to be specified. Canonical deconditioning is ground-attractive on the set of all states.

Fixing the ground already constrains every other state. We state this for a channel between two spaces, for use in the next subsection.

\begin{proposition}[A channel that carries the ground]
\label{prop:relaxation-carried}
Let $\mathcal T$ be a channel from the trace class of a Hilbert space $\mathcal H$ to that of $\mathcal H'$, and let $y\in\mathcal H$ and $y'\in\mathcal H'$ be unit vectors. Then $\mathcal T(P_y)=P_{y'}$ exactly when
\begin{equation}
 \delta_{y'}\bigl(\mathcal T(\sigma)\bigr)\le\delta_y(\sigma)\qquad\text{for every state $\sigma$ on $\mathcal H$.}
 \label{eq:relaxation-carried}
\end{equation}
\end{proposition}
\begin{proof}
If Eq.~\eqref{eq:relaxation-carried} holds, the state $\mathcal T(P_y)$ has no weight outside $y'$ and equals $P_{y'}$ by Proposition~\ref{prop:relaxation-pure-distance}. Conversely, let $F$ be the effect on $\mathcal H$ with $\Tr(F\sigma)=\Tr\bigl(Q_{y'}\mathcal T(\sigma)\bigr)$ for every trace-class $\sigma$. It exists because the right side is a positive linear functional of $\sigma$ bounded by the trace. If $\mathcal T(P_y)=P_{y'}$, then $\langle y,Fy\rangle=\Tr(Q_{y'}P_{y'})=0$. Since $F\ge0$, this gives $Fy=0$. Hence $F=Q_yFQ_y\le Q_y$ and $\Tr(F\sigma)\le\Tr(Q_y\sigma)$.
\end{proof}

A family that fixes the ground therefore never raises the weight outside it above its initial value, and under a semigroup that fixes the ground this weight is nonincreasing in time. The proposition uses that a channel preserves the trace. A conditional step, normalized after a selected outcome, can fix the ground and increase the weight outside it. With $K=\tfrac12P_y+Q_y$, the step $\sigma\mapsto K\sigma K/\Tr(K\sigma K)$ fixes $P_y$ and takes every state with weight $\delta$ outside the ground to one with weight $4\delta/(1+3\delta)$, because $\Tr(K\sigma K)=(1+3\delta)/4$ and $\Tr(Q_yK\sigma K)=\delta$. This is larger than $\delta$ for $0<\delta<1$ and is zero at $\delta=0$.

The step amplifies a weight that is already present. It does not take the ground anywhere, and no outcome of an instrument whose channel fixes the ground can. Let the operations $\mathcal O_k$ of the outcomes sum to a channel $\mathcal T$ with $\mathcal T(P_y)=P_y$. The positive operators $\mathcal O_k(P_y)$ then sum to the rank-one projection $P_y$, so each of them is a multiple of it. Every outcome that occurs with positive probability on the ground leaves the state $P_y$. A determination of Section~\ref{sec:determination} that takes $y$ to a different ray can therefore be an outcome only of an instrument whose channel does not fix the ground.

Fixing the ground does not make that weight decrease. If a unitary evolution fixes the ground, $U_tP_yU_t^*=P_y$, then $U_t^*Q_yU_t=Q_y$ and
\begin{equation}
 \delta_y\bigl(U_t\sigma U_t^*\bigr)=\Tr\bigl(U_t^*Q_yU_t\,\sigma\bigr)=\delta_y(\sigma)
 \label{eq:relaxation-unitary}
\end{equation}
for every state. The weight outside the ground never changes, so such an evolution is ground-attractive only on the basin $\{P_y\}$. A unitary can carry one chosen vector to $y$. It cannot do so while leaving $y$ in place. A ground-attractive evolution with a larger basin is therefore not unitary on the appearance space. It can still be the reduced description of a unitary evolution on a larger space.

Canonical deconditioning has such a description at each time. Let the space of an environment be $\mathcal H_{\mathrm e}=\C\oplus y^\perp$, with vectors written $(c,k)$, and define
\begin{equation}
 \iota_th=L_yh\otimes(1,0)+\sqrt{1-s^2}\;y\otimes(0,Q_yh),\qquad h\in\mathcal H_\Phi .
 \label{eq:relaxation-dilation}
\end{equation}
The two terms are orthogonal and $\|\iota_th\|^2=\|L_yh\|^2+(1-s^2)\|Q_yh\|^2=\|h\|^2$, so $\iota_t$ is an isometry into $\mathcal H_\Phi\otimes\mathcal H_{\mathrm e}$. The partial trace over the environment removes the cross terms between its two summands and gives
\begin{equation}
 \Tr_{\mathrm e}\bigl(\iota_t\sigma\iota_t^*\bigr)=L_y\sigma L_y+(1-s^2)\Tr(Q_y\sigma)P_y=\mathcal D_t^y(\sigma).
 \label{eq:relaxation-dilation-trace}
\end{equation}
Any unitary $U$ on $\mathcal H_\Phi\otimes\mathcal H_{\mathrm e}$ with $U\bigl(h\otimes(1,0)\bigr)=\iota_th$ then gives $\mathcal D_t^y(\sigma)$ as the reduced state of $U\bigl(\sigma\otimes P_{(1,0)}\bigr)U^*$. Such unitaries exist because the orthogonal complements of $\mathcal H_\Phi\otimes(1,0)$ and of $\Ran\iota_t$ have the same dimension. In this description the environment starts in $(1,0)$. The amplitude removed from $y^\perp$ is copied into the second summand of the environment while the system is placed in $y$, and afterwards the environment is found outside $(1,0)$ with probability $(1-s^2)\delta_y(\sigma)$. The construction shows that an environment of this kind is possible. It does not identify one in an observer.

The unitary above depends on $t$. On a finite-dimensional larger space, the weight outside the ground under one unitary group $e^{-itG}$ is a finite sum of terms $c_{ab}e^{-i(g_a-g_b)t}$ over the eigenvalues $g_a$ of $G$. Such a sum returns arbitrarily close to its initial value at arbitrarily late times. It cannot equal $e^{-\gamma t}\delta_y(\sigma)$ for all $t\ge0$ unless $\delta_y(\sigma)=0$. Reproducing the whole semigroup with one unitary group therefore requires an infinite-dimensional larger space, even when the appearance space has finite dimension. We do not construct one.

\subsection{Relaxation across many moments}
\label{subsec:relaxation-moments}
The momentary description lets the appearance space and its ground change from one moment to the next. Let $\mathcal H_{\Phi,n}$ be the appearance space at moment $n$, with ground vector $y_n$ and state $\sigma_n$, and let $\mathcal T_n$ be the channel with $\sigma_{n+1}=\mathcal T_n(\sigma_n)$. Put $\delta_n=\delta_{y_n}(\sigma_n)$. The spaces may be identified with one another, but nothing below requires it. Two conditions can be stated without choosing canonical deconditioning. The steps \emph{carry the ground} when
\begin{equation}
 \mathcal T_n(P_{y_n})=P_{y_{n+1}}\qquad\text{for every }n,
 \label{eq:relaxation-ground-carried}
\end{equation}
and the given sequence of states \emph{contracts with factors} $q_n\in[0,1]$ when
\begin{equation}
 \delta_{n+1}\le q_n\,\delta_n\qquad\text{for every }n .
 \label{eq:relaxation-contraction}
\end{equation}
The first is a property of the channels. The second concerns the particular states $\sigma_n$.

\begin{proposition}[Stabilization across moments]
\label{prop:relaxation-moments}
\begin{enumerate}
\item[(a)] If the steps carry the ground and $\sigma_n=P_{y_n}$ at some moment $n$, then $\sigma_m=P_{y_m}$ at every later moment $m$.
\item[(b)] If the sequence contracts with factors $q_n$, then
\begin{equation}
 \delta_n\le\Bigl(\prod_{k=0}^{n-1}q_k\Bigr)\delta_0,
 \qquad \|\sigma_n-P_{y_n}\|_1\le2\sqrt{\delta_n}.
 \label{eq:relaxation-product}
\end{equation}
If the products tend to zero, then $\|\sigma_n-P_{y_n}\|_1\to0$.
\item[(c)] Step $n$ carries the ground exactly when $\delta_{y_{n+1}}(\mathcal T_n(\sigma))\le\delta_{y_n}(\sigma)$ for every state $\sigma$ on $\mathcal H_{\Phi,n}$. Thus, when the steps carry the ground, every sequence of states contracts with the factors $q_n=1$.
\end{enumerate}
\end{proposition}
\begin{proof}
Part (a) is Eq.~\eqref{eq:relaxation-ground-carried} applied repeatedly. Iterating Eq.~\eqref{eq:relaxation-contraction} gives the first bound of (b), and Proposition~\ref{prop:relaxation-pure-distance} in each space gives the second. Part (c) is Proposition~\ref{prop:relaxation-carried} applied to the step.
\end{proof}

Part (b) does not use Eq.~\eqref{eq:relaxation-ground-carried}, and it allows steps with $q_n=1$ at which nothing is gained. By (c), carrying the ground is the same as contraction of every state with the factor one. It excludes any increase of the weight outside the ground, and it does not ensure a decrease: a unitary step that carries the ground leaves that weight unchanged for every state, by the computation of Eq.~\eqref{eq:relaxation-unitary}. By (b), factors whose products tend to zero guarantee stabilization. They are not necessary when the factors are bounds that hold for every state. On $\Span\{g,a,b\}$ with ground $g$, let one step be the channel with Kraus operators $P_g+P_b$ and $\ket{g}\!\bra{a}$, which resets $a$ to $g$ and leaves $b$ alone, and let the next step do the same with $a$ and $b$ exchanged. Both carry the ground. Each leaves one excited state unchanged, so neither contracts every state by a factor below one. Together they send every state to the ground. For an example with exact factors, let step $n$ be canonical deconditioning toward $y_n$ at rate $\gamma_n$ for a duration $\tau_n$, followed by an isometry from $\mathcal H_{\Phi,n}$ into $\mathcal H_{\Phi,n+1}$ that takes $y_n$ to $y_{n+1}$. These steps carry the ground, and by Eq.~\eqref{eq:relaxation-decay} every sequence of states satisfies Eq.~\eqref{eq:relaxation-contraction} with equality for $q_n=e^{-\gamma_n\tau_n}$. The products tend to zero exactly when $\sum_n\gamma_n\tau_n$ diverges. Count the steps from one. If the $n$th has $\gamma\tau=1/n$, every state approaches the ground. If it has $\gamma\tau=1/n^2$, the weight outside the ground decreases at every step and tends to $e^{-\pi^2/6}\delta_0$, about $0.19\,\delta_0$.

We say that the sequence \emph{is at the ground} at moment $n$ when $\sigma_n=P_{y_n}$, and that it \emph{stabilizes} when $\|\sigma_n-P_{y_n}\|_1\to0$. By (a), being at the ground persists when the steps carry it. The two notions differ. Under the deconditioning steps just described, a sequence that starts off the ground is not at the ground at any finite moment, and it stabilizes exactly when $\sum_n\gamma_n\tau_n$ diverges. In the cognitive model the first is an event of exact realization and the second is its asymptotic stabilization.

\subsection{What is returned to the physical side after relaxation?}
\label{subsec:relaxation-return}
Let $V:S\to\mathcal H_\Phi$ be an isometry with range $D$, as in Proposition~\ref{prop:two-moment-deconditioning}, and write $\Pi_D=VV^*$ and $Q_D=I-\Pi_D$. For an appearance state $\sigma$ the operator $V^*\sigma V$ on $S$ is positive with trace $\Tr(\Pi_D\sigma)=1-p_\perp(\sigma)$, in the notation of Eq.~\eqref{eq:matched-defect-outcome}. When this trace is positive, the \emph{returned state} is
\begin{equation}
 \sigma_S=\frac{V^*\sigma V}{\Tr(\Pi_D\sigma)},\qquad
 V\sigma_SV^*=\frac{\Pi_D\sigma\Pi_D}{\Tr(\Pi_D\sigma)} .
 \label{eq:relaxation-return}
\end{equation}
The second expression is $\sigma$ conditioned on passing the projective test $\Pi_D$, whose complement is the outcome $\perp$ of Section~\ref{sec:polar-matched-observables}. For a pure state $P_\phi$ the returned state is the ray of $V^*\phi$. It is defined exactly when $\phi$ does not lie in $D^\perp$. The map $\sigma\mapsto V^*\sigma V$ is completely positive and does not increase the trace. Any measurement with the effect $\Pi_D$ gives its outcome with probability $\Tr(\Pi_D\sigma)$. The state after that outcome depends on the instrument. It is $\sigma_S$ for an instrument whose operation for this outcome is $\sigma\mapsto V^*\sigma V$. An instrument that prepares a fixed state $\rho_0$ on $S$ after the same outcome has the operation $\sigma\mapsto\Tr(\Pi_D\sigma)\rho_0$. It gives the same probabilities and leaves $\rho_0$. The return is therefore a choice of operation. A later measurement on $S$ distinguishes two such instruments exactly when one of the effects available to it has different expectations in the states they leave.

Suppose $y\in D$, with partner $p=V^*y$. Equation~\eqref{eq:moment-excess-decay} gives, for every appearance state $\sigma$,
\begin{equation}
 \Tr\bigl(\Pi_D\,\mathcal D_t^y(\sigma)\bigr)=1-e^{-\gamma t}\,\Tr(Q_D\sigma).
 \label{eq:relaxation-paired-weight}
\end{equation}
This is positive for $t>0$, including for a state that began entirely outside $D$. Since $V^*P_yV=P_p$ and compression by $V$ does not increase the trace norm, Eq.~\eqref{eq:relaxation-distance} bounds the distance of $V^*\mathcal D_t^y(\sigma)V$ from $P_p$. Dividing a positive operator of trace $w\le1$ by $w$ changes it by $1-w$ in trace norm, and Eq.~\eqref{eq:relaxation-paired-weight} gives $w$. Hence
\begin{equation}
 \bigl\|\bigl(\mathcal D_t^y(\sigma)\bigr)_S-P_p\bigr\|_1
 \le2e^{-\gamma t/2}\sqrt{\delta_y(\sigma)}+e^{-\gamma t}\,\Tr(Q_D\sigma).
 \label{eq:relaxation-return-bound}
\end{equation}
The returned state converges to the partner of the ground from every initial appearance state. When $V$ is the polar isometry of $A$ and the weak bridge of Section~\ref{sec:ground-correspondence} holds, that partner is the physical vacuum $x$ exactly under strong polar compatibility.

A second return uses only the two reference lines. With $u_0=\ket{y}\!\bra{x}$ of Eq.~\eqref{eq:restricted-ground-orientation}, the operation $\sigma\mapsto u_0^*\sigma u_0=\langle y,\sigma y\rangle P_x$ has the effect $P_y$ and leaves $P_x$ whenever its outcome occurs. Equation~\eqref{eq:relaxation-decay} gives
\begin{equation}
 u_0^*\,\mathcal D_t^y(\sigma)\,u_0=\bigl(1-e^{-\gamma t}\delta_y(\sigma)\bigr)P_x .
 \label{eq:relaxation-restricted-return}
\end{equation}
The coefficient is positive for every $t>0$ and tends to one. After normalization this return is the physical vacuum at every positive time, by construction. The return through $V$ approaches the partner $p$, which is the vacuum only under strong compatibility.

On the unequal-cell path of Section~\ref{sec:ground-correspondence}, start from the fine site $e_2$ considered earlier in this section, for which $\Tr(Q_D\sigma)=3/4$ and $\delta_y(\sigma)=4/5$. At half relaxation, $s^2=1/2$, the state lies in $D$ with probability $5/8$ and passes the fine ground test with probability $3/5$. The state returned through $V$ converges to $(1,2)/\sqrt5$, which passes the coarse ground test with probability $9/10$. The restricted return gives the coarse ground itself.
